% !TEX root = ./../main.tex

\section{Related work}
\label{sec:related_work}
This section discusses existing diffusion-based generative models for videos including long video generation techniques.

\subsection{Video diffusion models}
\label{subsec:vdm}
Diffusion models, originally developed for high-quality image synthesis, have become a prominent approach in video generation~\citep{chen2023videocrafter1,ho2022video,singer2022make,zhou2022magic,wang2023modelscope}.
VDM~\citep{ho2022video} modifies the structure of U-Net~\citep{ronneberger2015unet} and proposes a 3D U-Net architecture to incorporate temporal information for denoising. 
On the other hand, Make-A-Video~\citep{singer2022make} employs a 1D temporal convolution layer following a 2D spatial convolutional layer to approximate 3D convolution.
This design enables the model to capture visual-textual relationships by training spatial layers with image-text pairs before incorporating temporal context through 1D temporal layers. 
Recently, \citep{peebles2023dit} introduce a transformer architecture,  known as DiT, for diffusion models.
Additionally, several open-sourced text-to-video models have emerged~\citep{wang2023modelscope,chen2023videocrafter1,wang2023lavie,chen2024videocrafter2}, trained on large-scale text-image and text-video datasets.


\subsection{Long video generation}
\label{subsec:long_vdm}

Long video generation approaches typically involve training models to predict future frames sequentially~\citep{voleti2022mcvd,harvey2022flexible,blattmann2023align, chen2023seine}. or generate a set of frames in a hierarchical manner~\citep{he2022latent,yin2023nuwa}.
For instance, Video LDM~\citep{blattmann2023align} and MCVD~\citep{voleti2022mcvd} employ autoregressive techniques to sequentially predict frames given several preceding ones, while FDM~\citep{harvey2022flexible} and SEINE~\citep{chen2023seine} generalize masked learning strategies for both prediction and interpolation.
Autoregressive methods are capable of producing indefinitely long videos in theory, but they often suffer from quality degradation due to error accumulation and limited temporal consistency across frames.
Alternatively, NUWA-XL~\citep{yin2023nuwa} adopts a hierarchical approach, where a global diffusion model generates sparse key frames with local diffusion models filling in frames using the key frames as references. 
However, this hierarchical setup requires batch processing, making it unsuitable for generating infinitely long videos. 

There are a few training-free long video generation techniques.
Gen-L-Video~\citep{wang2023genl} treats a video as overlapped short clips and introduces temporal co-denoising, which averages multiple predictions for one frame.
FreeNoise~\citep{qiu2023freenoise} employs window-based attention fusion to sidestep the limited attention scope issue and proposes local noise shuffle units for the initialization of long video.
FreeNoise requires memory proportional to the video length for the computation of cross, limiting its scalability for generating infinitely long videos.

\subsection{Diffusion models with latents of different noise levels}
Recent studies have adopted diffusion models for sequence generation by leveraging a sliding window approach with temporally varying noise levels~\citep{zhang2024tedi,ruhe2024rolling}.
These methods train diffusion models from scratch to accommodate latents with different noise levels, addressing tasks such as motion generation~\citep{zhang2024tedi} and video prediction~\citep{ruhe2024rolling}. 
However, training diffusion models from scratch introduces significant computational costs, especially for text-to-video generation tasks. 
In contrast, our approach is a training-free inference technique based on the standard diffusion models, trained on latents with uniform noise, for sequence generation within the sliding window framework. 
While \citep{ruhe2024rolling} is implemented with a nested loop to deal with two different axes corresponding to video frame index and diffusion time step, FIFO-Diffusion combines these two dimensions using a 1D queue, improving efficiency with a single loop.
